\section{Datasets and Lineage}

This project uses a layered data pipeline that starts from merged raw records and ends in a model-ready defensive-back (DB) table. Table~\ref{tab:source-lineage} documents the source-to-model lineage.

\begin{table}[htbp]
\centering
\small
\setlength{\tabcolsep}{5pt}
\caption{Source-to-model lineage for the DB modeling dataset.}
\label{tab:source-lineage}
\begin{tabular}{>{\raggedright\arraybackslash}p{0.20\linewidth} >{\raggedright\arraybackslash}p{0.33\linewidth} >{\raggedright\arraybackslash}p{0.40\linewidth}}
\hline
\textbf{Layer} & \textbf{Primary files} & \textbf{Role in pipeline} \\
\hline
Raw / merged base data & \path|all_data.csv|; \path|NFL_data/combine_with_college_stats.csv|; \path|extension_and_data/*.csv| & Core player-level records combining NFL season logs, combine measurements, and college statistics used for feature derivation and heuristic scoring. \\
Heuristic sweep outputs & \path|outputs/pipeline/sweeps/db_heuristic_grid_search/*| (especially \path|best_heuristic_scored_players.csv|) & Contains per-player/per-season heuristic outputs, including \path|NFL_production_value|, after evaluating candidate scoring configurations. \\
Modeling table & \path|artifacts/db_model_preprocessed.csv| & Final supervised-learning table with target, split labels, standardized predictors (\path|_z|), and missingness indicators (\path|_missing|). \\
\hline
\end{tabular}
\end{table}

\subsection{How the Data Was Collected}

The raw inputs were assembled from three collection paths before preprocessing:
\begin{itemize}
    \item \textbf{Combine data:} provided directly to the project team as source files.
    \item \textbf{NFL data:} scraped from ESPN using a non-public API.
    \item \textbf{College football data:} scraped from Pro Sports Reference using a Chrome extension workflow.
\end{itemize}
These sources were then merged into the unified base tables referenced in Table~\ref{tab:source-lineage}.

\subsection{Competition Dataset Fields vs. External Additions}

\textbf{Competition/base fields.} The competition-ready modeling inputs are restricted to pre-draft information families:
(1) combine/anthropometric fields (e.g., \texttt{40yd}, \texttt{Vertical}, \texttt{Bench}, \texttt{Broad Jump}, \texttt{3Cone}, \texttt{Shuttle}, \texttt{Ht}, \texttt{Wt}) and
(2) college production fields (e.g., \texttt{college\_games}, \texttt{college\_interceptions}, \texttt{college\_passes\_defended}, \texttt{college\_combined\_tackles}, \texttt{college\_tfl}, \texttt{college\_sacks}).

\textbf{External additions.} Additional artifacts are introduced by the project pipeline rather than the original merged files:
\begin{itemize}
    \item \texttt{NFL\_production\_value}: heuristic target value generated during sweep experiments.
    \item \texttt{dataset\_split}: train/validation/test assignment.
    \item Derived model features from preprocessing: standardized columns (suffix \texttt{\_z}) and explicit missingness indicators (suffix \texttt{\_missing}).
\end{itemize}

\subsection{Missing-Value Strategy and Outlier Treatment}

Preprocessing is implemented in \texttt{src/data/preprocess\_db\_model.py}. The pipeline:
\begin{enumerate}
    \item filters to DB-family positions using \texttt{Pos} and \texttt{college\_pos};
    \item collapses season-level rows to one row per player (summing \texttt{NFL\_production\_value} across retained seasons and taking first values for covariates);
    \item coerces mixed-format measurement fields (e.g., height and broad jump) into numeric inches;
    \item applies deterministic median imputation for each base feature;
    \item appends binary missingness flags for each imputed feature;
    \item creates z-score standardized versions of each feature for modeling.
\end{enumerate}

Outlier handling is intentionally conservative: there is no explicit winsorization, clipping, or row deletion for high/low numeric values in this preprocessing stage. Instead, values are retained and mapped into standardized scale through z-scores, preserving extreme observations while reducing scale imbalance.

\subsection{Target Construction and First-Contract Window Rationale}

The supervised target is \texttt{NFL\_production\_value}, sourced from the heuristic sweep output table and carried into the final modeling artifact. In feature preprocessing, \texttt{pipeline/features/preprocessing.py} enforces an early-career restriction to the first four NFL seasons using
\[
\texttt{career\_year} = \texttt{season\_year} - \texttt{combine\_year} + 1 \le 4.
\]
This first-contract window aligns target construction with the rookie-contract evaluation horizon and focuses model learning on the period most relevant to recruiting and early career projection.
